\section{Описательный анализ данных}
\label{sec:dataset_description}

Опубликованный набор данных iPinYou включает несколько типов логов, отражающих ключевые этапы RTB-процесса: журналы ставок в аукционах (bidding log), журналы показов/кликов/конверсий (impression / click / conversion log), а также вспомогательные таблицы (категории рекламодателей, профили пользователей, соответствия регионов и городов). Ниже приведено описание основных сущностей и полей.

\subsection{Журнал ставок}
Каждая строка журнала содержит сведения из трёх измерений: (1) пользователь, (2) рекламный слот, (3) параметры ставки и рекламной кампании. Формат данных представлен в Таблице \ref{tab:bidding_log_format}.

Дадим некоторым полям более подробное описания:

\begin{itemize}
    \item \textbf{Timestamp} -- момент времени, когда запрос ставки (bid request) поступает на DSP в формате yyyyMMddHHmmssSSS.
    \item \textbf{iPinYou ID} -- идентификатор пользователя (cookie), установленный iPinYou.
    \item \textbf{User Agent} -- информация об операционной системе клиента, типе браузера, модели устройства.
    \item \textbf{IP address} -- первые три байта IP-адреса пользователя. Последний байт удаляется в целях защиты приватности.
    \item \textbf{Ad exchange} -- идентификатор рекламной биржи, с которой пришёл данный аукцион. Возможные значения: 1--6, соответствующие Tanx (Alibaba), Adx (Google DoubleClick AdX), Tencent (Tencent), Baidu (Baidu), Youku (Youku), Amx (Google Mobile) соответственно.
    \item \textbf{URL, Anonymous URL ID} -- URL страницы, на которой должен быть показан рекламный слот. Издатель может запретить передавать этот URL в DSP, в таком случае используется поле \textbf{Anonymous URL ID}. 
    \item \textbf{Ad slot ID} -- уникальный идентификатор рекламного слота (места размещения) на странице. Одна страница может содержать несколько рекламных слотов.
    \item \textbf{Ad slot width}, \textbf{Ad slot height} -- ширина и высота рекламного слота. Поля могут быть массивами, что означает допустимость нескольких размеров креатива для данного слота.
    \item \textbf{Ad slot visibility} -- видимость слота: на верхней части страницы (FirstView), или ниже (SecondView-TenthView), неизвестно (Na).
    \item \textbf{Ad slot format} -- формат слота: фиксированное место (Fixed), всплывающее окно (Pop), неизвестно (Na).
    \item \textbf{Ad slot floor price} -- минимальная цена, при которой издатель допускает победу DSP в аукционе. Если ни одна ставка не превышает минимальную, победителя нет. Перед публикацией применяется линейное масштабирование значения. Единица измерения: юань за CPM.
    \item \textbf{Creative ID} -- идентификатор рекламного креатива рекламодателя, по которому DSP делает ставку.
    \item \textbf{Bidding price} -- ставка DSP в аукционе. Перед публикацией применяется линейное масштабирование значения. Единица измерения: юань за CPM.
    \item \textbf{User Profile IDs} -- тэги пользователя, однако в bidding log эти значения всегда пусты по соображениям безопасности. Поле представляет собой массив числовых идентификаторов категорий.
\end{itemize}

\begin{table}[h!]
\small
  \begin{center}
    \caption{Формат данных bidding log. Столбцы, отмеченные символом \textasteriskcentered, содержат значения, которые были захэшированы или модифицированы перед публикацией лога.}
    \label{tab:bidding_log_format}
    \renewcommand{\arraystretch}{1.15}
    \begin{tabular}{c|c|l|p{10.5cm}}
      \textbf{N} & \textbf{*} & \textbf{Column} & \textbf{Example} \\
      \hline
      1  & \textasteriskcentered & Bid ID & c0550000008e5a94ac18823d6f275121 \\
      2  &  & Timestamp & 20130218134701883 \\
      3  & \textasteriskcentered & iPinYou ID & 35605620124122340227135 \\
      4  &  & User-Agent &
      Mozilla/5.0 (Windows NT 5.1) AppleWebKit/535.11 (KHTML, like Gecko)\\
      &&& Chrome/17.0.963.84 Safari/535.11 SE 2.X MetaSr 1.0 \\
      5  & \textasteriskcentered & IP & 119.163.222.\textasteriskcentered \\
      6  &  & Region ID & 146 \\
      7  &  & City ID & 147 \\
      8  &  & Ad Exchange & 2 \\
      9  & \textasteriskcentered & Domain & e80f4ec7f5bfbc9ca416a8c01cd1a049 \\
      10 & \textasteriskcentered & URL & hz55b000008e5a94ac18823d6f275121 \\
      11 &  & Anonymous URL & null \\
      12 &  & Ad Slot ID & 9737269023493 \\
      13 &  & Ad Slot Width & 300 \\
      14 &  & Ad Slot Height & 250 \\
      15 &  & Ad Slot Visibility & FirstView \\
      16 &  & Ad Slot Format & Na \\
      17 & \textasteriskcentered & Ad Slot Floor Price & 0 \\
      18 &  & Creative ID & f80f4ec7f5bfbc9ca416a8c01cd1a049 \\
      19 & \textasteriskcentered & Bidding Price & 573 \\
      20 &  & Advertiser ID & 2259 \\
      21 & \textasteriskcentered & User Profile IDs & null \\
    \end{tabular}
  \end{center}
\end{table}

\subsection{Журнал показов / кликов / конверсий}
В журналах показов, кликов и конверсий поля совпадают с журналом ставок. Помимо полей из журнала ставок дополнительно присутствуют следующие поля:

\begin{itemize}
    \item \textbf{Log type} -- тип события: 1 (показ), 2 (клик), 3 (конверсия).
    \item \textbf{Paying price} -- фактическая цена, которую платит победившая DSP на бирже. Перед публикацией, аналогично ставкам, применяется линейное масштабирование. Единица измерения: юань за CPM.
    \item \textbf{Landing page URL} -- URL страницы, на которую переходит пользователь при клике по креативу.
\end{itemize}

Формат которых описан в Таблице \ref{tab:imp_clck_conv_log_format}:

\begin{table}[h!]
\small
  \begin{center}
    \caption{Формат данных impression/click/conversion log. Столбцы, отмеченные символом \textasteriskcentered, содержат значения, которые были захэшированы или модифицированы перед публикацией лога.}
    \label{tab:imp_clck_conv_log_format}
    \renewcommand{\arraystretch}{1.15}
    \begin{tabular}{c|c|l|p{10.5cm}}
      \textbf{N} & \textbf{*} & \textbf{Column} & \textbf{Example} \\
      \hline
      1  &  & Bid ID & 01530000008a77e7ac18823f5a4f5121 \\
      2  &  & Timestamp & 20130218134701883 \\
      3  &  & Log Type & 1 \\
      4  & \textasteriskcentered & iPinYou ID & 35605620124122340227135 \\
      5  &  & User-Agent & Mozilla/5.0 (compatible; MSIE 9.0; Windows NT 6.1; WOW64; Trident/5.0) \\
      6  & \textasteriskcentered & IP & 118.81.189. \\
      7  & \textasteriskcentered & Region ID & 15 \\
      8  &  & City ID & 16 \\
      9  &  & Ad Exchange & 2 \\
      10 & \textasteriskcentered & Domain & e80f4ec7f5bfbc9ca416a8c01cd1a049 \\
      11 & \textasteriskcentered & URL & hz55b000008e5a94ac18823d6f275121 \\
      12 &  & Anonymous URL & null \\
      13 &  & Ad Slot ID & 21476898764813 \\
      14 &  & Ad Slot Width & 300 \\
      15 &  & Ad Slot Height & 250 \\
      16 &  & Ad Slot Visibility & SecondView \\
      17 &  & Ad Slot Format & Fixed \\
      18 & \textasteriskcentered & Ad Slot Floor Price & 0 \\
      19 &  & Creative ID & e39e178fdf366606f8cab791ee56bcd \\
      20 & \textasteriskcentered & Bidding Price & 753 \\
      21 & \textasteriskcentered & Paying Price & 15 \\
      22 & \textasteriskcentered & Landing Page URL & a8be178fdf366606f8cab791ee56bcd \\
      23 &  & Advertiser ID & 3358 \\
      24 & \textasteriskcentered & User Profile IDs & 123,5678,3456 \\
    \end{tabular}
  \end{center}
\end{table}

\subsection{Категории рекламодателей}
Для идентификации рекламодателя используется поле \textbf{Advertiser ID}. Полный перечень категорий рекламодателей приведён в Таблице \ref{tab:advertiser_category}.

\begin{table}[h!]
\small
  \begin{center}
    \caption{Категории рекламодателей.}
    \label{tab:advertiser_category}
    \renewcommand{\arraystretch}{1.15}
    \begin{tabular}{l|l}
      \textbf{Advertiser Key} & \textbf{Industrial Category} \\
      \hline
      1458 & Chinese vertical e-commerce \\
      3358 & Software \\
      3386 & International e-commerce \\
      3427 & Oil \\
      3476 & Tire \\
      2259 & Milk powder \\
      2261 & Telecom \\
      2821 & Footwear \\
      2997 & Mobile e-commerce app install \\
    \end{tabular}
  \end{center}
\end{table}


\subsection{Профиль пользователя}
iPinYou разработала классификацию аудитории для описания характеристик пользователей по нескольким аспектам: демография, география, долгосрочные интересы (long-term) и покупательские намерения (in-market).

Данные профилей представляют собой таблицу соответствий между идентификатором категории и её названием.

\subsection{Таблица регионов и городов}
iPinYou предоставляет два файла: (1) соответствие \textbf{region id} $\rightarrow$ \textbf{region name} и (2) соответствие \textbf{city id} $\rightarrow$ \textbf{city name}. Все значения регионов и городов относятся к географическим локациям материкового Китая.
 
\subsection{Предварительная обработка данных}
\label{subsec:data_preprocessing}

Отметим, перед публикацией данных iPinYou исключает из bidding log площадки с подозрительной активностью. Для обнаружения спама используются специализированные алгоритмы фильтрации \cite{ipinyou-contest-database}.

В ходе валидации качества данных при помощи скрипта были выявлены и устранены несколько типов аномалий и несоответствий. Во-первых, в таблице bidding log в поле user agent встречались некорректные значения: одиночный символ '':'', использовавшийся вместо пропуска, а также строки, содержащие IP-адрес. Для таких записей были заданы пустые значения user agent. Во-вторых, было обнаружено, что для части идентификаторов городов, присутствующих в bidding log, отсутствует соответствие в таблице городов (не задано текстовое наименование города) -- подобные случаи будут обработаны базой данных во время операций \textbf{JOIN}.

Кроме того, в данных встречались записи, где идентификатор рекламной биржи отсутствовал, несмотря на строгий формат возможных значений 1-6. Для обеспечения целостности данных и корректности группировок при последующем анализе было введено дополнительное значение категории -- 0, соответствующее случаю Unknown. Наконец, при подготовке файлов к загрузке в СУБД были выявлены различия в кодировке переводов строк (Windows CRLF vs Unix LF): данное несоответствие было устранено на этапе предобработки скриптом для обеспечения корректного парсинга строк.

Для повышения ценности данных на основе поля user agent с помощью отдельной программы-парсера были сформированы два дополнительных признака: browser (тип браузера) и platform (устройство/платформа). Полученные признаки далее используются при сегментации трафика и анализе зависимостей метрик эффективности от технических характеристик пользователя.

